\itemhead{p.37 练习 · 第2题（教材p.37）}{XJ-B1-C02-PR-P037-Y1148}
\label{XJ-B1-C02-PR-P037-Y1148}
{\footnotesize 湘教版教材p.37 / PDF第44页 / p.37 练习 / 第2题}
\field{原题}{回答下列问题：
\subq{1}若 $a>b$，且 $c>d$，能否判断 $a-c$ 与 $b-d$ 的大小？举例说明。
\subq{2}若 $a>b$，且 $c<d$，能否判断 $a+c$ 与 $b+d$ 的大小？举例说明。
\subq{3}若 $a>b$，且 $c>d$，能否判断 $ac$ 与 $bd$ 的大小？举例说明。
\subq{4}若 $a>b,c<d$，且 $c\ne0,d\ne0$，能否判断 $\dfrac ac$ 与 $\dfrac bd$ 的大小？举例说明。}
\field{参考答案}{四问均不能。各给出两组 $(a,b,c,d)$：（1）$(3,1,2,1)$、$(2,1,4,1)$；（2）$(3,1,0,1)$、$(2,1,0,3)$；（3）$(2,1,2,1)$、$(-1,-2,-1,-2)$；（4）$(2,1,1,2)$、$(3,1,-2,-1)$。每组均满足该小问题设，所得比较方向相反。}
{\small \field{关联知识}{I02 不等式基本性质与符号条件（教材帮p.45、46 / PDF54、55）}
\field{方法与易错点}{每一步明确所乘除因子的正负与非零，或作差求证。 易错：未知符号就直接乘除；将 $c^2$ 当作必为正数而忘记 $c=0$。}
\field{答案依据}{答案册相关题组：PDF第11页（册内页10）；按当前题库原题逐项重算，不按旧版题号机械对套。}}
